%==========================================================
\subsection{Derivations and error propagation}
\label{sec:deriv}
This section sets out the short derivations behind the closed-form layers, so that each can be audited line by line, and the elementary error analysis that explains which deviations in the verification were to be expected.

Assume a sharp wetting front at depth $L_f$ with suction head $\psi_f$ and a saturated transmission zone of conductivity $K_s$. Darcy's law across the wetted zone gives the infiltration rate
\begin{equation}
f=\frac{\mathrm dF}{\mathrm dt}=K_s\Bigl(1+\frac{\psi_f\Delta\theta}{F}\Bigr),\qquad F=L_f\Delta\theta,
\label{eq:GAode}
\end{equation}
and separating variables, $\int_0^{F}\bigl(1+\psi_f\Delta\theta/F'\bigr)^{-1}\mathrm dF'=K_st$, yields \cref{eq:GA} \citep{GreenAmpt1911}. Two limits are used as sanity checks in the code: for small $t$ the logarithm expands as $F-\psi_f\Delta\theta\ln(1+F/\psi_f\Delta\theta)\approx F^2/(2\psi_f\Delta\theta)$, so $F\approx\sqrt{2K_s\psi_f\Delta\theta\,t}$ (sorptivity-dominated, $F\propto t^{1/2}$); for large $t$, $F\to K_st$ (gravity-dominated). The Kostiakov exponent $A$ fitted to $F(\tau)$ over 10--600\,min is therefore expected to lie between $0.5$ and $1$ (the fitted exponents have not been tabulated for the ten soils; the interface displays them with the fit $R^2$). Because \cref{eq:GAode} is monotone, the Newton iteration of the code converges from the starting value $F_0=\sqrt{2K_s\psi_f\Delta\theta t}+K_st$.

Let the advance front be at $x(t)$ and the inflow be $Q$. Conservation of volume over $[0,t]$ states that the inflow equals the volume stored on the surface plus the volume infiltrated:
\begin{equation}
Qt=\sigma_yA_0\,x+\sigma_z\,Z(t)\,w\,x,
\label{eq:vbderiv}
\end{equation}
where $A_0$ is the flow cross-section at the inlet, $\sigma_y$ the surface-storage shape factor (the ratio of the mean to the inlet cross-section along the wetted length) and $\sigma_z$ the sub-surface shape factor. For a power-law advance $x=pt^{r}$ and Kostiakov intake, the infiltrated volume per unit width is $\int_0^x Z(t-t_x)\,\mathrm dx$, which evaluates to $\sigma_zZ(t)x$ with $\sigma_z$ as in \cref{eq:vb} \citep{Elliott1982,Walker1987}. Solving \cref{eq:vbderiv} for $x$ gives \cref{eq:vb}, and the exponent $r$ is found from the two measured (or predicted) times $t_{L/2}$ and $t_L$ by $r=\ln 2/\ln(t_L/t_{L/2})$, since $x\propto t^{r}$ implies $L/(L/2)=(t_L/t_{L/2})^{r}$. Because $\sigma_z$ depends on $r$ and $r$ on the times, the solution is a fixed point, found by iteration; the code reports non-convergence rather than returning a number.

Let $V_{\mathrm{inf}}=V_{\mathrm{st}}+V_{\mathrm{dp}}$, which holds by construction since $\min(Z,Z_n)+\max(Z-Z_n,0)=Z$ pointwise. If the infiltrated volume does not exceed the applied volume, $V_{\mathrm{ro}}=V_{\mathrm{app}}-V_{\mathrm{inf}}\ge0$, so
\begin{equation}
E_a+\mathrm{DP}+\mathrm{RO}=\frac{V_{\mathrm{st}}+V_{\mathrm{dp}}+V_{\mathrm{ro}}}{V_{\mathrm{app}}}=\frac{V_{\mathrm{inf}}+V_{\mathrm{app}}-V_{\mathrm{inf}}}{V_{\mathrm{app}}}=1.
\label{eq:effclose}
\end{equation}
The identity fails only when the numerical $V_{\mathrm{inf}}$ exceeds $V_{\mathrm{app}}$, which signals a violated volume balance and is flagged in the interface. The same argument gives $E_a+\mathrm{DP}+\mathrm{EV}=1$ for the drip layer, where evaporation takes the place of runoff.

For $n$ catch-can depths $a_k$ with mean $\bar a$, the Christiansen coefficient $\mathrm{CU}=1-\sum|a_k-\bar a|/(n\bar a)$ equals one minus the relative mean absolute deviation \citep{Christiansen1942}. For normally distributed depths the mean absolute deviation is $\sigma\sqrt{2/\pi}$, so $\mathrm{CU}\approx1-0.80\,\sigma/\bar a$; this relation allows the coefficient of variation to be inferred from a reported CU and was used to translate the ``CU $\ge85$\,\%'' check into a coefficient of variation of at most $0.19$. For the drip layer, \cref{eq:Eu} of \citet{KellerKarmeli1974} treats the lowest quartile of emitters; the factor $1.27$ is the standardised mean of the lowest quarter of a normal distribution, so that $1-1.27\,c_v$ is the ratio of the low-quarter mean to the mean discharge, and division by $\sqrt{n_e}$ accounts for the averaging of $n_e$ emitters per plant; $n_e=1$ gives the most conservative estimate.

For a lateral of $N$ emitters at spacing $s_e$, the flow in the segment that feeds emitter $j$ (counting from the inlet) is $\sum_{i\ge j}q_i$, and the head at emitter $j$ is the inlet head less the accumulated losses. With uniform emitters and the exponent $x_e$, the total head loss for a pipe of length $L$ with outlet factor $F_c$ is approximated by Christiansen's factor
\begin{equation}
h_f\approx F_c\,\frac{10.67\,L\,Q_0^{1.852}}{C^{1.852}D^{4.871}},\qquad F_c=\frac1{2.852}+\frac1{2N}+\frac{\sqrt{1.852-1}}{6N^2},
\label{eq:Fc}
\end{equation}
which tends to $1/2.852=0.351$ for large $N$ and is the textbook shortcut used as an independent plausibility check on the segment-by-segment iteration. Because $q_j\propto h_j^{x_e}$ with $x_e<1$, the loop \emph{head $\to$ flow $\to$ loss $\to$ head} is a contraction when the relative loss $h_f/H_0$ is small, with contraction ratio of the order $x_e\,h_f/H_0$; four iterations then reduce the initial error by a factor of order $(x_eh_f/H_0)^4$, which is of order $10^{-4}$ (about $2\times10^{-4}$ for $x_e=0.6$) for $h_f/H_0\le0.2$ and much larger for $h_f/H_0\ge0.5$ (where the ratio is of order $0.3$). This simple analysis predicts the qualitative result of \cref{tab:drip} before any numerical test, and it explains why the failures occur only in poorly designed laterals.

Engine~1 returns products of ratios, \cref{eq:E1E,eq:E1T,eq:E1ET,eq:E1I,eq:E1Y}. If each ratio carries an independent relative error $\delta_i$, the relative error of the product is $\delta\approx\sqrt{\sum\delta_i^2}$; with the three to five factors per flux and ratio errors of a few per cent this gives a few per cent, so random error cannot explain the $+53$\,mm bias reported in Section~\ref{sec:e1}, which is consistent with a \emph{systematic} baseline offset ($\delta_b=(720-673)/673=7\,\%$ for ET and $(282-228)/228=24\,\%$ for irrigation). The additive treatment of drainage and leaching propagates absolute errors, $\sigma_D^2=\sum\sigma_i^2$, and because the mean drainage is small (35\,mm) a modest absolute error is a large relative one, which explains the normalised RMSE of 48.7\,\% despite $R^2=0.82$. Finally, the closure window of \cref{eq:E1close} bounds the storage drift to $-15\le\Delta S\le30$\,mm\,yr$^{-1}$, which over a 20-year simulation could accumulate to a few hundred millimetres in the extreme. The window is not a rare safeguard: in the held-out test of Section~\ref{sec:e1} it is active in 680 of the 1\,000 cells (482 at the upper and 198 at the lower bound of ET), so in that experiment Engine~1's ET is mostly set by the closure $P+I-D$ and not by the ET ratios, which is consistent with its weak ET skill ($R^2=-0.11$ after re-anchoring, $\rho=0.84$). The user-facing balance reports $\Delta S$, so the clipping is visible.

Rank correlation is invariant to monotone transformations, whereas $R^2$ and RMSE are not. A systematic multiplicative or additive offset shifts RMSE and $R^2$ but leaves Spearman's $\rho$ unchanged; this is why the raw irrigation $R^2$ of 0.18 becomes 0.79 after re-anchoring while $\rho$ is 0.96 in both cases. A screening emulator is judged by whether it orders the options correctly, a budgeting model by whether it reproduces the magnitude, and the verification therefore reports both families of metric \citep{Moriasi2007,Bennett2013}.


